\documentclass[12pt]{article}
\usepackage[utf8]{inputenc}
\usepackage[T1]{fontenc}
\usepackage{hyperref}
\title{Xav\'ante Language Resources for K-12 Curriculum Development}
\author{Compiled by Claude Sonnet 5.5 (claude-sonnet-5-5) for Andr\'as Kornai}
\date{}
\begin{document}
\maketitle

\section{Introduction}
Xav\'ante (A'uw\~e; ISO 639-3: xav; Glottocode: xava1240) is a Je-family language of Mato Grosso, Brazil. Wikipedia reports about 9,600 speakers, of whom about 7,000 are monolingual (2006 figure). It is taught in some schools. Do not confuse it with Xerente (Akw\~e), a separate language.

\section{Primary Educational Resources}
Welch's 2023 book on A'uw\~e life cycles says its native terms follow the orthography currently used by A'uw\~e teachers at the Pimentel Barbosa Municipal School and the Etênhiritip\'a State Fundamental Education School. I saw no URL for the book itself. The searches turned up no standalone textbooks or primers with URLs. A Brazilian thesis on Xavante age classes and schooling gives pedagogical context: \url{https://site.ucdb.br//public/md-dissertacoes/7922-classes-de-idade-xavante-e-a-educacao-escolar-uma-contribuicao-pedagogica.pdf}

\section{Linguistic and Reference Materials}
\begin{itemize}
\item \emph{A'uw\~e Mreme: Dicion\'ario Preliminar da L\'ingua Xavante} (2020), Huff, Lapierre, Carrick and Xavante co-authors: \url{https://escholarship.org/uc/item/8j764483}
\item Webonary Xavante dictionary, English introduction (it mentions a grammar and some instruction in the language): \url{https://www.webonary.org/xavante/visao-geral/english-introduction/?lang=en}
\item Xavante language (Wikipedia): \url{https://en.wikipedia.org/wiki/Xavante_language}
\item Article in Entreletras (UFNT): \url{https://periodicos.ufnt.edu.br/index.php/entreletras/article/download/994/532/3591}
\end{itemize}

\section{Religious and Community Resources}
None were found in the searches.

\section{Digital and Audio Resources}
The Webonary site above is an online lexicon. The searches turned up no separate audio or video collections.

\section{Curriculum Development Context}
Xavante is spoken by a large community that has bilingual indigenous schools and a developing orthography. Writing is being revitalized by A'uw\~e researchers, as in the thesis \emph{Wamreme: a revitaliza\c{c}\~ao da escrita da l\'ingua a'uw\~e uptabi}: \url{https://repositorio.bc.ufg.br/tede/items/1a2be540-ea69-4670-936e-0e9f9f49267d}. The preliminary dictionary and the Webonary lexicon are the natural base for curriculum vocabulary.
\end{document}